\section{Strategy 2 Rigor Completion}
\label{sec:strategy2-rigor-completion}

This section supplies the two pieces that must not be hidden by the concise
reader presentation: the complete four-label T3-like endpoint audit and the
explicit Vd1--supercritical axis replacement.  The former is long enough that
its complete branch algebra is maintained as an electronic supplement in the
proof package.  The supplement is part of the proof: the paths and Git blob
identifiers below pin the exact source objects at this manuscript revision.

\subsection{The incorporated \(407X\) endpoint supplement}

The support-isolated T3-like proof is distributed as follows.
\begin{center}
\scriptsize
\renewcommand{\arraystretch}{1.08}
\begin{tabular}{p{.42\textwidth}p{.18\textwidth}p{.30\textwidth}}
\toprule
proof-package file&Git blob&role\\
\midrule
\texttt{4073\_boundary\_loss\_framework.md}&
\texttt{f3f0395748f5}&exact residual reduction\\
\texttt{4074\_L\_Full\_branch.md}&
\texttt{880b59ec8d3c}&$(L,\mathrm{Full})$ branch\\
\texttt{4075\_Tminus\_low\_lower\_branch\_obligations.md}&
\texttt{ebe2b65f8840}&all first-$T_-$ branches\\
\texttt{4078\_left\_L\_family\_completion.md}&
\texttt{f1323325b069}&remaining first-$L$ branches\\
\texttt{4079\_first\_Full\_branch.md}&
\texttt{c66eea05277d}&first-Full exclusion\\
\texttt{407a\_left\_Thigh\_branch\_completion.md}&
\texttt{aa5cb1a6dc63}&first-$T_+^{\rm hi}$ branches\\
\texttt{407c\_rigor\_completion\_details.md}&
\texttt{c80243f67124}&all omitted analytic inequalities\\
\texttt{407d\_rigor\_final\_assembly.md}&
\texttt{a6ef81f787c4}&exhaustive reassembly\\
\bottomrule
\end{tabular}
\end{center}
All paths are relative to
\begin{verbatim}
proof/4XXX_CE1CE2/40XX_Nplus0/
407X_T3_like_no_Vd1Vd2/
\end{verbatim}
The exact capped-map theorem used there is
Proposition~\ref{prop:capped-demand-map}.  Its four labels are precisely
\[
 \mathrm{Full},\qquad L,\qquad T_-,\qquad T_+^{\rm hi};
\]
there is no fifth formal sheet.

For reference, we reproduce the three non-obvious high-sheet estimates that
connect the branch files.

\begin{lemma}[High-left envelope]
\label{lem:paper-high-left-envelope}
Put
\[
 \rho=\sqrt{r^2-r+1},\qquad
 m=\sqrt{\beta^2-\beta+1},
\]
where
\[
 0<r<1,\qquad
 \beta\ge\max\left\{r,\frac12,\frac{1-r^2}{1+2r}\right\},
\]
and assume the selected Cell-$T$ condition
\[
 r\ge(1-\beta)(r+\beta)^2.
\]
Define
\[
 y_*=\frac{r}{1-r}
 \left(\frac{m}{r+\beta}-\frac12-rac{1-r}{1+\rho}\right),
 \qquad b=\frac{\beta m}{r+\beta}.
\]
Whenever the companion endpoint $S$ satisfies $S<1/2$, the exact center
identity gives $y<y_*$, and
\[
 y_*<\frac1{10},
 \qquad
 3y_*\le1-b.
\]
\end{lemma}

\begin{proof}
The inequality $y<y_*$ is the rearrangement of $S<1/2$.  The function
$m/(r+\beta)$ decreases in $\beta$, so the first bound is checked at
\[
 \beta_0(r)=\max\left\{r,\frac12,\frac{1-r^2}{1+2r}\right\}.
\]
Let $r_0=(\sqrt3-1)/2$.  If $0<r\le r_0$, then
$m/(r+\beta)\le1$ and $(1-r)/(1+\rho)\ge1/3$, whence
$y_*<1/10$.  If $r_0\le r\le1/2$, use
\[
 \frac{m}{r+\beta}\le\frac{\sqrt3}{1+2r},
 \qquad
 \frac{1-r}{1+\rho}\ge
 2-\sqrt3+\frac{1/2-r}{2}.
\]
The latter follows after squaring from
\[
 -\frac{(r-1)(2r-1)}{16}
 \left(2r^2+(8\sqrt3-17)r+56-32\sqrt3\right)\ge0.
\]
The resulting difference from $1/10$ has numerator
\[
 P(r)=20r^3+(40\sqrt3-96)r^2+(40\sqrt3-57)r-2,
\]
which is increasing and negative at $1/2$; the denominator has the opposite
sign.  Finally, if $1/2\le r<1$, then
\[
 y_*\le\frac{\rho+1-3r}{2(1+\rho)}<\frac1{10},
\]
because $15r-4\rho-4>0$ and its derivative is positive.

For the second inequality put $E_r(\beta)=1-b-3y_*$.  Exact differentiation
gives
\[
 \frac{\partial E_r}{\partial\beta}
 =\frac{N_r(\beta)}{2(\beta+r)^2(1-r)m},
\]
where
\[
 N_r=2\beta^3r-2\beta^3+4\beta^2r^2-5\beta^2r+\beta^2
      -9\beta r^2+5r^2+4r.
\]
Moreover
\[
 \partial_\beta N_r
 =-6\beta^2(1-r)+2\beta(1-r)(1-4r)-9r^2<0.
\]
Thus $E_r$ has no interior minimum.  At $\beta=1$,
\[
 E_r(1)=\frac{r(6r-\rho+5)}{2(1+r)(1+\rho)}>0.
\]
At the lower endpoint, the three ranges
$r\le r_0$, $r_0\le r\le1/2$, and $r\ge1/2$ give respectively
\[
 \frac{r(7+\rho-2r-8r\rho-14r^2-2r^2\rho)}
 {2(1-r)(1+2r)(1+\rho)}>0,
\]
$E_r>1/8$, and
\[
 E_r(r)=\frac{10r-r^2-2\rho-2}{2(1+\rho)}>0.
\]
Hence $E_r\ge0$.
\end{proof}

\begin{lemma}[Right \(T_-\) bounds]
\label{lem:paper-right-tminus}
Let $C=1-y$, $q=tC$, and suppose the right endpoint has the selected
$T_-$ label.  If $b_1$ is its capped output, then
\[
 q+b_1=\sqrt{t^2-t+1},
\]
and
\[
 b_1\le\frac{\sqrt{13}-1}{6},
 \qquad
 q\le\tau<\frac{93}{200},
\]
where $\tau\in(0,1/2)$ is the unique root of
\[
 x^4-3x^3+3x^2-3x+1=0.
\]
\end{lemma}

\begin{proof}
The branch condition $b_1\le q$ gives
$b_1\le\frac12\sqrt{t^2-t+1}$ and
$\sqrt{t^2-t+1}\le2t$.  Hence
$3t^2+t-1\ge0$ and
$t\ge(\sqrt{13}-1)/6$.  Since $y<1/10$ and $q<1/2$, one has $t<5/9$; the
half-radical is at most $(\sqrt{13}-1)/6$ on this interval.

The selected Cell-$T$ inequality is
\[
 -tC^2+\sqrt{t^2-t+1}\,C-(1-t)(t^2-t+1)\ge0.
\]
It first forces $t<1/2$.  Its larger $C$-root is
$\sqrt{t^2-t+1}(1-t)/t$, so
\[
 q\le\min\left\{t,(1-t)\sqrt{t^2-t+1}\right\}.
\]
The first function increases and the second decreases.  Their intersection is
the displayed quartic root; direct evaluation at $93/200$ is negative.
\end{proof}

\begin{lemma}[Analytic high-left threshold]
\label{lem:paper-high-left-threshold}
Under the hypotheses of Lemma~\ref{lem:paper-high-left-envelope}, set
\[
 S_0=1-\frac{m}{r+\beta}+\frac{1-r}{1+\rho},
 \qquad B_5=\frac{\beta m}{r+\beta}.
\]
Then
\[
 S_0<\frac{93}{200}
 \quad\Longrightarrow\quad
 B_5<\frac{5657}{10000}.
\]
\end{lemma}

\begin{proof}
The selected Cell-$T$ condition factors as
\[
 r-(1-\beta)(r+\beta)^2
 =(r+\beta-1)(\beta^2+r\beta-r).
\]
The second factor is nonnegative, so $s=r+\beta\ge1$.  Since $r\le\beta$,
\[
 \rho^2-m^2=(r-\beta)(s-1)\le0.
\]
Define
\[
 \overline S(\beta)=1-m+\frac\beta{1+m}.
\]
Then
\[
 S_0\ge\overline S(\beta),
\]
because the difference is
\[
 (s-1)\left(\frac ms-\frac1{1+m}\right)\ge0;
\]
here $s\le2\beta\le m(1+m)$, the latter following from
\[
 m^2-(3\beta-\beta^2-1)^2
 =-\beta(\beta-1)(\beta^2-5\beta+5)\ge0.
\]
Put
\[
 B_0=\frac{5657}{10000},\qquad T=\frac{93}{200},
 \qquad \beta_*=\frac{161}{250}.
\]
If $B_5\ge B_0$, then $\beta m\ge B_0$.  The function $\beta m$ is increasing,
and
\[
 B_0^2-\beta_*^2m(\beta_*)^2
 =\frac{22782769}{62500000000}>0,
\]
so $\beta>\beta_*$.  On $[\beta_*,1]$ define
\[
 F(\beta)=\frac{107}{200}-Tm-(1-\beta)^2.
\]
Since $F''=-3T/(4m^3)-2<0$, $F$ is concave.  Its endpoint values are positive:
$F(1)=7/100$, and, with
$a_*=51033/125000$,
\[
 a_*^2-T^2m(\beta_*)^2
 =\frac{1693881}{62500000000}>0.
\]
Finally
\[
 (\overline S(\beta)-T)(1+m)=F(\beta)>0.
\]
Thus $B_5\ge B_0$ forces $S_0>T$; the contrapositive proves the lemma.
\end{proof}

\begin{theorem}[Complete support-isolated T3-like endpoint certificate]
\label{thm:paper-complete-t3-endpoint}
For every support-isolated state of
Definition~\ref{def:body-t3-endpoint-state},
\[
 \widehat g_{c_5^*}(1-a_5^*)+
 \widehat g_{c_1^*}(1-a_1^*)<1.
\]
\end{theorem}

\begin{proof}
The exact residual formulas are those displayed in
Lemma~\ref{lem:407-four-label-loss}.  If $a_1^*+a_5^*>1$, the two diagonal caps
already give the strict loss.  In the hard region, the exact capped-map theorem
leaves the four first-coordinate labels listed above.

For first label $L$, the low-root separation, the CE2-overlap comparison, and
the center-transfer lemma are proved in the incorporated files
\texttt{4074}, \texttt{4078}, and the first section of \texttt{407c}.  They
cover right labels Full, $L$, $T_-$, and $T_+^{\rm hi}$.  For first label
$T_-$, the exact root-order test and all non-overlap, overlap, and right-high
subcases are proved in \texttt{4075}; in every genuine case either
$B_5<a_1^*$ or the two low caps sum to less than one.  The first-Full label is
excluded in \texttt{4079} separately for the non-overlap input
$1-\alpha$ and the CE2 overlap input $1-T$.

For first label $T_+^{\rm hi}$, the selected Cell-$T$ condition is the
hypothesis of Lemma~\ref{lem:paper-high-left-envelope}.  Right Full and right
high are excluded by $S>3y$ and $A_C>3y$; right $L$ follows from
$e(y)<3y\le1-B_5$; and the two right-$T_-$ inputs are handled by
Lemmas~\ref{lem:paper-right-tminus} and
\ref{lem:paper-high-left-threshold}, together with the exact comparison
\[
 \frac{5657}{10000}<\frac{7-\sqrt{13}}6.
\]
These are precisely the branches proved in \texttt{407a} and \texttt{407c}.
The four labels are exhaustive and the exact final assembly is
\texttt{407d}.  Thus every branch has endpoint-cap sum strictly below one.
\end{proof}

\subsection{The explicit Vd1--supercritical replacement}

\begin{proposition}[Vd1--supercritical axis replacement]
\label{prop:paper-vd1-pair-replacement}
Assume the CE2 exact-$M_0$ branch, with a unique supercritical row at $V_1$ and
an adjacent Vd1 row at $V_0$ containing $M_1$.  Assume every other vertex row
is nonsupercritical Vd0, no additional role has positive support on the two
relevant radial handoffs, and the center has no boundary trace on the shared
edge $e_{0,1}$.  Then the two special rows can be replaced by two open
nonsupercritical Vd0 roles preserving all boundary and radial demands used by
the all-Vd0 boundary-loss theorem.
\end{proposition}

\begin{proof}
Use the Vd corner normal form with $t>0$, $d=\sqrt{t^2+t+1}$ and exact Vd1
boundary reaches $a,b$:
\[
 x-(t+1)y\le a,\qquad
 ty-(t+1)x\le tb,\qquad
 tx+y\le d-a-tb.
\]
Its own-radial reach and neighboring-arm endpoints are
\[
 c=\frac{d-a-tb}{t+1},\qquad
 \lambda=\frac{t(1-b)}{t+1},\qquad
 \mu=\frac{d-a-tb-1}{t}.
\]
Because the original open Vd1 role contains $M_1$,
\[
 b>\frac{t-1}{2t},\qquad
 a+tb<d-1-\frac t2.
\]
If $t<1$, then
\[
 d-a-tb-t>1-\frac t2>\frac1{t+1}>\frac{1-a}{t+1},
\]
which would create positive support on the other adjacent arm.  Hence
$t\ge1$.

We need four strict margins:
\begin{equation}
 a+c<1,\qquad a<\lambda\le\frac12,\qquad
 b<\frac12,\qquad \mu<1-a.
 \label{eq:paper-replacement-margins}
\end{equation}
For the first, maximize $a+c$ on the closed midpoint face
$a=d-1-t/2-tb$.  Then
\[
 a+c-1=-\frac{2bt^2+2bt-2dt-2d+t^2+4t+2}{2(t+1)}.
\]
Using $b\ge(t-1)/(2t)$, the numerator is at least
$2t^2+4t+1-2(t+1)d$, which is positive because
\[
 (2t^2+4t+1)^2-4(t+1)^2d^2
 =4t^3+4t^2-4t-3>0
\]
for $t\ge1$.  The original midpoint inequalities are strict.
The bound $\lambda\le1/2$ follows from the lower bound for $b$.  The closed
upper bound for $a$, subtracted from $\lambda$, is minimized at
$b=(t-1)/(2t)$ and equals $(t+1)-d>0$, so $a<\lambda$.
Furthermore
\[
 b<\frac{d-1-t/2}{t}<\frac12
\]
because $d<t+1$.  Finally
\[
 t(\mu+a-1)=d-(t+1)-tb+(t-1)a
 \le t(d-t-1)<0.
\]
This proves \eqref{eq:paper-replacement-margins}.

Let $(a_1,b_1,c_1)$ be the actual reaches of the unique supercritical row.
The center-free shared edge forces
\[
 a_1\ge1-b>\frac12,
\]
and radial coverage before the Vd1 interval forces $c_1\ge\lambda$.  We use the
following admissibility consequence.

\emph{Half-square lemma.}
If $(x,y,z)\in\mathcal A$, $x\ge1/2$, and $0\le z\le1/2$, then
\[
 y\le1-z.
\]
Indeed, suppose $y>1-z$.  Then $x+y>1$, so the selected supercritical cell
applies.  In the ordered half $x\le y$, its necessary polynomial is
\[
 F(x,y,z)=(x^2-1)z^2+(2xy^2+y)z+y^4-y^2\le0.
\]
It is nondecreasing in $x$.  At $x=1/2$ its $y$-derivative is
$4y^3-2y+2yz+z$, which is increasing for $y\ge1-z$ and is positive at that
endpoint.  Hence
\[
 F(x,y,z)>F\left(\frac12,1-z,z\right)
 =\frac{z^2(2z-5)(2z-1)}4\ge0,
\]
a contradiction.  In the reflected half $y<x$, use $F(y,x,z)$.  Both
arguments exceed $r=1-z$, and on that quadrant the two partial derivatives are
nonnegative; the second is bounded below by
\[
 4r^2z+z+4r^3-2r=2-5z+4z^2>0.
\]
Thus
\[
 F(y,x,z)>F(r,r,z)=z(1-2z)\ge0,
\]
again a contradiction.

A supercritical row has selected own-radial coordinate at most $1/2$, so the
lemma applied to $(a_1,b_1,c_1)$ gives
\[
 b_1\le1-c_1\le1-\lambda.
\]
Together with $a<\lambda$ this gives
\begin{equation}
 a+b_1<1.
 \label{eq:paper-a-bi-margin}
\end{equation}

Let $d_1^C$ be the center reach on the rescued arm, measured from $O$, and put
$c_1^{\rm req}=1-d_1^C$.  Only the center, the supercritical row, and the Vd1
row have positive intervals on this arm.  Open coverage at the center handoff
therefore gives
\[
 c_1^{\rm req}<\max\{c_1,\mu\}
 <\max\{1-b_1,1-a\}.
 \label{eq:paper-radial-bridge-margin}
\]
Equality would leave the common open exit/entry point uncovered.

Choose
\[
 a<p_1<p_2<1-b_1
\]
within the strict margins so that
\[
 p_1<\frac12,\qquad 1-p_1>c,
 \qquad \max\{p_2,1-p_2\}>c_1^{\rm req}.
\]
This is possible by choosing $p_1$ close to $a$ and choosing $p_2$ close to
$a$ when the right side of \eqref{eq:paper-radial-bridge-margin} is supplied by
$1-a$, or close to $1-b_1$ when it is supplied by $1-b_1$.

In the local metric define
\[
 \Delta_p^-=\conv\{(0,1-p),(1,1-p),(0,-p)\}
 \quad(p\le1/2),
\]
\[
 \Delta_p^+=\conv\{(p,0),(p,1),(p-1,0)\}
 \quad(p\ge1/2).
\]
These are unit equilateral triangles.  Their incident reaches are $p,1-p$;
their own-radial reaches are respectively $1-p,p$; and they have no positive
adjacent-arm support.

Translate $\Delta_{p_1}^-$ by $(-\varepsilon,0)$.  Its relevant reaches become
\[
 (p_1-\varepsilon,1-p_1,1-p_1).
\]
For the second row translate $\Delta_{p_2}^-$ by $(-\varepsilon,0)$ if
$p_2<1/2$, and translate $\Delta_{p_2}^+$ by $(0,-\varepsilon)$ if
$p_2\ge1/2$.  Its relevant reaches are respectively
\[
 (p_2-\varepsilon,1-p_2,1-p_2)
\]
or
\[
 (p_2,1-p_2-\varepsilon,p_2).
\]
Choose $\varepsilon>0$ smaller than every strict boundary, radial, and
$p_2-p_1$ margin.  The first replacement still supplies at least $a$ and $c$;
the second supplies at least $b_1$ and $c_1^{\rm req}$.  The shared open edge
is covered because
\[
 (1-p_1)+(p_2-\varepsilon)>1.
\]
Both replacements contain their distinguished vertices in their interiors,
have zero adjacent support, and have boundary sum strictly below one.  Thus
they are open nonsupercritical Vd0 roles, and every demand used by the all-Vd0
boundary-loss theorem is preserved.
\end{proof}
